% Copyright 2020-2025 Robert Bosch GmbH

% Licensed under the Apache License, Version 2.0 (the "License");
% you may not use this file except in compliance with the License.
% You may obtain a copy of the License at

% http://www.apache.org/licenses/LICENSE-2.0

% Unless required by applicable law or agreed to in writing, software
% distributed under the License is distributed on an "AS IS" BASIS,
% WITHOUT WARRANTIES OR CONDITIONS OF ANY KIND, either express or implied.
% See the License for the specific language governing permissions and
% limitations under the License.

% --------------------------------------------------------------------------------------------------------------

\section{Overview}\label{sec:overview}
The \textbf{EventBusClient} provides a high-level API for interacting with a \textbf{RabbitMQ}-based event bus system.
It abstracts low-level AMQP details and supports dynamic loading of project-specific plugins for serializers, exchange handlers, and message structures.
This client is designed for modularity and ease of use across multiple projects.


\section{Features}\label{sec:features}
\begin{itemize}
  \item \textbf{Dynamic plugin loading}: Load custom serializers, exchange handlers, and message classes at runtime from a configurable plugins directory.
  \item \textbf{Thread-safe publishing}: Supports safe publishing from multiple threads with minimal locking overhead.
  \item \textbf{Automatic reconnect}: Automatically reconnects to \textbf{RabbitMQ} after connection loss (if enabled in configuration).
  \item \textbf{Pluggable architecture}: Supports both built-in and project-specific extensions.
  \item \textbf{Supports JSON, Protobuf, Pickle serializers}.
\end{itemize}

\section{Architecture}\label{sec:architecture}
The \textbf{EventBusClient} is built around a modular architecture that allows for easy extension and customization. It consists of the following components:
\begin{itemize}
  \item \textbf{Base Classes}: Define interfaces for exchange handlers, serializers, and messages.
  \item \textbf{Plugins Directory}: Contains custom implementations of the base classes.
  \item \textbf{Configuration File}: Specifies the plugins path, \textbf{RabbitMQ} connection details, and selected implementations.
  \item \textbf{EventBusClient Class}: The main class that interacts with \textbf{RabbitMQ} and manages plugins.
  \item \textbf{Message Classes}: Define the structure of messages sent and received over the event bus.
  \item \textbf{Exchange Handlers}: Handle the declaration and publishing of messages to \textbf{RabbitMQ} exchanges.
  \item \textbf{Serializers}: Convert messages to and from byte streams for transmission over \textbf{RabbitMQ}.
\end{itemize}
The client uses a combination of built-in plugins and user-defined plugins to provide flexibility in message handling and serialization.
\section{Installation}\label{sec:installation}
To install the \textbf{EventBusClient}, use pip to install the package from PyPI:
\begin{pythonlog}
pip install event-bus-client
\end{pythonlog}
Ensure you have \textbf{RabbitMQ} server running and accessible at the specified host and port in the configuration file.
\section{Dependencies}\label{sec:dependencies}

The \textbf{EventBusClient} requires the following dependencies:
\begin{itemize}
  \item \textbf{pika}: For synchronous \textbf{RabbitMQ} communication
  \item \textbf{aio-pika}: Specifically for asynchronous \textbf{RabbitMQ} communication
  \item \textbf{protobuf}: For Protocol Buffers serialization (if using \textbf{ProtobufSerializer})
  \item \textbf{jsonpickle}: For JSON serialization (if using \textbf{JSONSerializer})
\end{itemize}
Ensure these dependencies are installed in your Python environment. You can install them using pip:
\begin{pythonlog}
pip install pika protobuf jsonpickle
\end{pythonlog}
\section{Directory Structure}\label{sec:directory-structure}
\begin{pythonlog}
project/
|-- EventBusClient/
|   |-- event_bus_client.py
|   |-- connection.py
|   |-- plugin_loader.py
|   |-- publisher.py
|   |-- subscriber.py
|   |-- qlogger.py
|   |-- message/
|   |   |-- base_message.py
|   |-- exchange_handler/
|   |   |-- base.py
|   |   |-- xr_topic_handler.py
|   |-- serializer/
|   |   |-- base_serializer.py
|   |   |-- json_serializer.py
|   |   |-- protobuf_serializer.py
|   |   |-- pickle_serializer.py
|   |-- plugins/
|   |   |-- mcpi/
|   |   |   |-- serialize/
|   |   |   |-- message/
|   |   |   |-- exchange/
|   |-- config/
|   |   |-- config.jsonp
|-- app.py
\end{pythonlog}

\newpage
\section{Configuration}\label{sec:configuration}
The \plog{config.jsonp} file controls all aspects of the \textbf{EventBusClient}, including plugin paths, \textbf{RabbitMQ} connection, and selected implementations.

\begin{pythoncode}
{
  "plugins_path": "./plugins",
  "host": "localhost",
  "port": 5672,
  "serializer": "PickleSerializer",
  "exchange_handler": "XRTopicExchangeHandler",
  "message_class": "ListenerEventMsg",
  "threadsafe_publish": true,
  "auto_reconnect": true,
  "qos_prefetch": 10
}
\end{pythoncode}


\section{Usage}\label{sec:usage}

\subsection{Initialization}
Create an \textbf{EventBusClient} instance from configuration:
\begin{pythoncode}
from event_bus.client import EventBusClient

client = await EventBusClient.from_config("./config/config.jsonp")
\end{pythoncode}


\subsection{Publishing Messages}
Send a message using the configured serializer and exchange handler:
\begin{pythoncode}
msg = ListenerEventMsg(event="start_cycle", timestamp=1234567890)
await client.send("zone1.topic.start", msg)
\end{pythoncode}

Enable thread-safe publish from a non-async thread:
\begin{pythoncode}
client.send("zone1.topic.alert", msg, threadsafe=True)
\end{pythoncode}


\subsection{Subscribing to Messages}
Subscribe to a routing key and handle incoming messages:
\begin{pythoncode}
async def on_message(msg: ListenerEventMsg):
    print(f"Received event: {msg.event} at {msg.timestamp}")

await client.on("zone1.#", ListenerEventMsg, on_message)
\end{pythoncode}


\newpage
\section{Example: Producer and Consumer}
This example (from \plog{EventBusClient/test/test.py}) starts a producer and a consumer in separate processes:
\begin{pythoncode}
# Start consumer
consumer = Process(target=run_process, args=(consumer_process, config_path))
consumer.start()

# Start producer
producer = Process(target=run_process, args=(producer_process, config_path))
producer.start()
# Wait for both to finish
producer.join()
consumer.join()
\end{pythoncode}

\section{Example: Custom Plugin Structure}\label{sec:example:-custom-plugin-structure}
To create a custom plugin, follow these steps:
\begin{itemize}
  \item Create a directory for your plugin in the configured plugins path.
  \item Implement the required interfaces (e.g., BaseExchangeHandler, BaseMessage, BaseSerializer).
  \item Register your plugin in the \plog{config.jsonp} file.
  \item Ensure your plugin is discoverable by the \textbf{EventBusClient}.
  \item Place your plugin code in the plugins directory, following the structure:
  \item For example, if your plugin is named \plog{CustomPlugin}, the directory structure should look like this:
        \begin{pythonlog}
plugins/
|-- CustomPlugin/
    |-- __init__.py
    |-- custom_exchange_handler.py
    |-- custom_message.py
    |-- custom_serializer.py\end{pythonlog}
\end{itemize}
\section{Example: Built-in Plugins}
The \textbf{EventBusClient} comes with several built-in plugins that can be used directly or extended:
\begin{itemize}
  \item \textbf{XRTopicExchangeHandler}: Handles topic exchanges with \textbf{RabbitMQ}.
  \item \textbf{ListenerEventMsg}: A default message class for listener events.
  \item \textbf{PickleSerializer}: Serializes messages using Python's Pickle module.
  \item \textbf{JSONSerializer}: Serializes messages to JSON format.
  \item \textbf{ProtobufSerializer}: Serializes messages using Protocol Buffers.
\end{itemize}
These plugins can be used as-is or extended to create custom functionality.

\newpage
\section{Example: Built-in Configuration}
The built-in configuration file \plog{config.jsonp} provides a starting point for configuring the \textbf{EventBusClient}:
\begin{pythoncode}
{
  "plugins_path": "./plugins",
  "host": "localhost",
  "port": 5672,
  "serializer": "PickleSerializer",
  "exchange_handler": "XRTopicExchangeHandler",
  "message_class": "ListenerEventMsg",
  "threadsafe_publish": true,
  "auto_reconnect": true,
  "qos_prefetch": 10
}
\end{pythoncode}

This configuration specifies the plugins path, \textbf{RabbitMQ} connection details, serializer, exchange handler, message class, and other options.


\section{Example: Custom Exchange Handler}\label{sec:example:-custom-exchange-handler}
To create a custom exchange handler, implement the required interface and place it in the plugins directory. For example, a custom exchange handler might look like this:
\begin{pythoncode}
from event_bus.plugins import BaseExchangeHandler
class CustomExchangeHandler(BaseExchangeHandler):
    def declare_exchange(self, channel):
        # Custom exchange declaration logic
        pass

    def publish(self, channel, routing_key, message):
        # Custom publish logic
        pass
# Register the plugin in config.jsonp
{
  "plugins_path": "./plugins",
  "exchange_handler": "CustomExchangeHandler"
}
\end{pythoncode}

\section{Example: Custom Message Class}\label{sec:example:-custom-message-class}
To create a custom message class, define it in your project and register it in the configuration:
\begin{pythoncode}
from event_bus.plugins import BaseMessage
class CustomMessage(BaseMessage):
    def __init__(self, data):
        self.data = data

    @classmethod
    def from_data(cls, data):
        pass

    @abstractmethod
    def get_value(self):
        pass
# Register the message class in config.jsonp
{
  "plugins_path": "./plugins",
  "message_class": "CustomMessage"
}
\end{pythoncode}

\section{Example: Custom Serializer}\label{sec:example:-custom-serializer}
To create a custom serializer, implement the required interface and place it in the plugins directory. For example, a custom serializer might look like this:
\begin{pythoncode}
from event_bus.plugins import BaseSerializer
class CustomSerializer(BaseSerializer):
    def serialize(self, obj):
        # Custom serialization logic
        pass

    def deserialize(self, data):
        # Custom deserialization logic
        pass
# Register the plugin in config.jsonp
{
  "plugins_path": "./plugins",
  "serializer": "CustomSerializer"
}
\end{pythoncode}
